\chapter{Смотрим шоу: партитура}\label{ch:partitura}

\textbf{Что вы узнете:} как по партитуре понять, что делает каждый дрон в любую секунду и какая строка кода за это отвечает.

\section{Время одно на всё окно}

Ползунок времени внизу общий для карты, партитуры, карточек и кода. Передвиньте его, и всё окно покажет один и тот же момент. Открывать файл для этого урока лучше всего тот, что мы уже знаем:

\solfile{l07_square.py}

Времена шагов сложим подряд: калибровка 0--5\,с (строка 7), взлёт 5--10 (строка 8), переход в клетку 13: 10--11 (строка 9), квадрат 11--27 (строка 10), посадка 27--30 (строка 11).

\section{Подсветка в коде}

Когда вы двигаете время, в коде подсвечивается строка, которая выполняется сейчас. Над кодом написано: «В 8,0 с у дрона 1 идёт строка 8». Если выбран не один дрон, а все, будет «идут строки …». Так можно быстро найти, какая строка отвечает за то, что вы видите на карте.

Второй цвет подсветки отмечает последнюю команду цвета, которая уже сработала.

\section{Полосы дронов}

Партитура — это лента времени: по одной полосе на каждого дрона, время идёт слева направо.

\begin{itemize}
  \item \textbf{Цветная полоса} показывает цвет светодиода во времени. Когда светодиод выключен, полосы нет.
  \item \textbf{Значки над полосой}: закрашенный треугольник вверх — взлёт, треугольник вниз — посадка, закрашенный квадрат — калибровка, вертикальная черта — переход в клетку, горизонтальная — смена высоты или ожидание, кружок — круг.
  \item \textbf{Линейка}: чёрточка каждые 10 секунд, число каждые 30.
  \item \textbf{Циановая линия} — текущее время.
\end{itemize}

Если навести мышь на полосу, появится подсказка вида «Дрон 2, 12,0 с: …, строка 28»: что делает дрон в этот момент и в какой строке кода это написано.

\section{Как двигаться по времени}

\begin{itemize}
  \item Мышью: нажмите на партитуру и тяните по ней (ползунок снизу работает так же).
  \item Стрелками: \klavisha{←} и \klavisha{→} сдвигают время на секунду, с \klavisha{Shift} на десять секунд. Сначала нажмите на партитуру, чтобы она получила фокус.
  \item \klavisha{↑} и \klavisha{↓} выбирают соседнего дрона; выбранный остаётся яркой полосой, остальные тускнеют.
\end{itemize}

\section{Карточки, вкладки и окна}

\begin{itemize}
  \item \textbf{Карточки} в режиме «Код»: по одной на дрона, с номером, высотой, цветом и строкой, которая идёт сейчас. Если шага нет, написано «ждёт» (дрон ещё не начал) или «закончил».
  \item \textbf{Вкладки} \knopka{Дрон 1}, \knopka{Дрон 2}…: показывают код только одного дрона. Они есть, когда шоу написано блоками \kod{d = show.drone(N)} по одному на дрона.
  \item \textbf{Окна кода} в режиме «Карта»: когда партитура раскрыта, по бокам карты стоят окна с кодом каждого дрона; в них видна строка, которая идёт сейчас, и по две строки вокруг. Окно можно закрыть крестиком.
\end{itemize}

\begin{probuy}
  \item Какая строка файла \texttt{l07\_square.py} идёт на 8-й секунде?
  \item А на 20-й секунде?
  \item На какой секунде дрон приходит в клетку 52?
  \item Сколько секунд длится шоу?
  \item Откройте пример \knopka{Разминуться} (\texttt{quarter\_pair\_cross}). Выберите дрона 2 стрелкой \klavisha{↓} и найдите в партитуре его взлёт. Какого цвета у него светодиод?
\end{probuy}
\otvety{9}
